\section{Исследование и построение решения задачи}
\label{sec:Chapter3} \index{Chapter3}

\subsection{Конечные автоматы и возможные ошибки в них}

\subsubsection{Основные определения}

\textbf{Конечный автомат} --- шестёрка $(S, I, O, \delta, \lambda, s_0)$, где:

\begin{itemize}
    \item $S$ --- конечное множество состояний автомата;
    \item $I$ --- множество входных сигналов;
    \item $O$ --- множество выходных сигналов;
    \item $\delta : S \times I \rightarrow S$ --- функция переходов;
    \item $\lambda$ --- функция выходов (зависит от типа автомата);
    \item $s_0 \in S$ — начальное состояние.
\end{itemize}

Конечный автомат в RTL --- последовательностная логическая модель, описывающая поведение системы
через конечное множество дискретных \textit{состояний}, \textit{переходов} между этими состояниями
и соответствующих им \textit{выходов}. В любой момент времени автомат находится ровно в
одном из своих возможных состояний. Следующее состояние определяется \textit{функцией переходов},
которая зависит от \textit{текущего состояния} и \textit{входов}
автомата. В RTL-коде состояние хранится в выделенном регистре, сами состояния кодируются именованными
константами, а \textit{входами} автомата служат сигналы, управляющие условиями переходов ---
выражениями в конструкциях if, case и тернарных операторах.

По виду функции выходов различают два классических типа конечных автоматов:

\begin{itemize}
\item \textbf{Автомат Мура} --- выходы зависят только от текущего состояния. Выход определяется
содержимым регистра состояния и не меняется на текущем такте при изменении входов. Формально это означает,
что функция переходов имеет следующий вид:

$$
\lambda: S \rightarrow O
$$

\item \textbf{Автомат Мили} — выходы зависят от текущего состояния и от текущих значений входов.
Такая модель позволяет получать быстрый отклик на входные сигналы, но усложняет верификацию
из-за комбинационного пути от входа к выходу. Функция переходов имеет вид:

$$
\lambda: S \times I \rightarrow O
$$

\end{itemize}

В настоящей работе рассматриваются автоматы обоих типов. Поскольку для обнаружения рассматриваемых видов
ошибок (недостижимые состояния, мёртвые переходы, дедлоки, лайвлоки) функция выходов несущественна,
различие между Муром и Мили в анализе сводится к единообразной обработке.

Автомат можно представить в виде ориентированного графа $G = (V, E)$, где:

\begin{itemize}
    \item $V = S$ --- множество вершин (состояний)
    \item $E \subseteq S$ --- множество переходов
\end{itemize}

Каждому ребру $e = (s_i, s_j)$ сопоставляется булево условие перехода $cond_{ij}$, определённое над
входами автомата $I$.

\subsubsection{Конечные автоматы в языках описания аппаратуры}

В языке SystemVerilog конечный автомат обычно реализуется через комбинацию регистра состояния и
комбинационной или последовательностной логики, определяющей следующее состояние. Канонический шаблон
описания выглядит следующим образом: вводится тип для состояний (обычно с помощью typedef enum),
объявляется регистр состояния и описывается один или два always-блока для логики переходов.
В односоставном варианте (одноблочный стиль) вся логика помещается в always\_ff, и переходы оформляются
в виде оператора case. В двухсоставном варианте (двухблочный стиль) комбинационная логика вычисления
следующего состояния выносится в always\_comb, а сам регистр обновляется в отдельном always\_ff.

На практике встречаются также варианты, в которых регистр состояния реализуется как экземпляр
параметризованного модуля триггера (например, шаблон из ядра VeeR-EL2, где функция следующего состояния
вычисляется в одном процедурном блоке, а регистр инстанцируется как модуль rvdff). Встречаются автоматы,
в которых в одном always-блоке совмещаются логика сброса и логика переходов; автоматы с несколькими
регистрами состояния; автоматы, реализованные через if/else вместо case. Эта вариативность стилей
существенно затрудняет разработку универсального анализатора.

Дополнительную сложность вносят параметризованные модули: один и тот же фрагмент исходного кода может
порождать конечные автоматы с различным набором состояний и переходов в зависимости от значений параметров,
подставленных при инстанцировании. Корректный анализ таких автоматов требует работы на уровне netlist,
полученного после параметризации (элаборации).

\subsubsection{Классификация типичных ошибок в конечных автоматах}

В реализации конечных автоматов в аппаратуре встречаются следующие типичные ошибки:

\textbf{Недостижимое состояние} --- состояние, в которое автомат не может попасть из начального
состояния ни при каком допустимом входном воздействии. Причины: опечатки в именах
состояний, неверные условия переходов, либо неактуальный код, оставшийся после рефакторинга.

\textbf{Мёртвый переход} --- переход, который никогда не происходит по каким-либо причинам (например, 
условие которого никогда не выполняется ни при каких входных
воздействиях).

Мёртвые переходы можно разделить на 2 типа:

\begin{itemize}
    \item Переходы, условия которых никогда не выполняются. Причинами таковых могут быть несовместные
    условия (например, выражение \texttt{ctl \& !ctl} всегда ложно), сигналы, жёстко привязанные к
    константам параметризацией или оптимизацией, а также логические ошибки в выражениях условий.
    \item Переходы, происходящие из недостижимых состояний
\end{itemize}

В рамках анализа будет рассматриваться первый тип мёртвых переходов, вторые же будут обнаруживаться
автоматически во время анализа недостижимых состояний.

\textbf{Дедлок (блокировка)} --- состояние, из которого не существует ни одного перехода в
другое состояние. В зависимости от проекта дедлок может быть как ошибкой, так и намеренно
спроектированным завершающим состоянием (например, после обнаружения необратимой аварии).

\textbf{Лайвлок («живая» блокировка)} --- группа состояний, переходы между которыми образуют
замкнутый цикл, но ни одно из них не имеет перехода в состояние вне этой группы.
Как и в случае дедлока, лайвлок может быть
намеренным (например, финальная фаза многофазного автомата, которая не возвращается
в начальное состояние).

\begin{table}[h!]
    \centering
    \begin{tabular}{cc}
    \makecell[b]{
        \includegraphics[width=0.26\linewidth]{unreachable.png} \\
        Недостижимое состояние (помечено красным)
    }
 &
    \makecell[b]{
        \includegraphics[width=0.15\linewidth]{dead_transition.png} \\
        Мёртвый переход (помечен красным)
    }
 \\
    \makecell[b]{
        \includegraphics[width=0.13\linewidth]{deadlock.png} \\
        Дедлок (помечен фиолетовым)
    }
 &
    \makecell[b]{
        \includegraphics[width=0.20\linewidth]{livelock.png} \\
        Лайвлок (помечен синим)
    }
 \\
    \end{tabular}
    \caption{Типы ошибок в конечных автоматах}
    \label{tab:fsm_errors}
\end{table}

Автоматический анализатор должен уметь обнаруживать все четыре вида ошибок. При этом нужно
понимать, что не каждый выявленный паттерн является реальной ошибкой — многие лайвлоки и дедлоки
соответствуют намеренной логике. Это требует ручной проверки и классификации результатов пользователем.

\subsubsection{Сильно связные компоненты и их роль в анализе}

\textbf{Сильно связной компонентой (ССК)} графа называют максимальный его подграф, в котором из
каждой вершины достижима любая другая.

\textbf{Главной группой} называется ССК, содержащая начальное состояние. Эта компонента
представляет собой «рабочую часть» автомата --- множество состояний, между которыми автомат
может курсировать в обычной работе.

\textbf{Терминальной сильно связной компонентой (ТССК)} называется ССК, отличная от главной
группы и не имеющая исходящих рёбер в вершины вне компоненты. ТССК соответствует ситуации,
когда автомат, попав в одно из её состояний, больше никогда не сможет покинуть эту группу.
Именно ТССК лежат в основе формального определения дедлоков и лайвлоков: ТССК из одной вершины
есть дедлок, ТССК из нескольких вершин --- лайвлок.

Для нахождения ССК применяется классический алгоритм Тарьяна, имеющий временную сложность
$O(|V| + |E|)$, где $|V|$ и $|E|$ --- число вершин и рёбер графа соответственно. Это позволяет
выполнять анализ даже на больших автоматах за приемлемое время.

\subsection{Инструмент SVAN и библиотека slang}

SVAN [3] --- это разрабатываемая в ИСП РАН система статического анализа для языка SystemVerilog.
Проект направлен на создание промышленного анализатора, способного работать с реальными
индустриальными RTL-проектами, включая проекты, критичные с точки зрения безопасности. 

Ключевым компонентом SVAN на этапе разбора исходного кода является открытая библиотека
slang [15] --- инструментарий для работы с языком SystemVerilog на уровне АСД. Библиотека
предоставляет полноценный лексический и синтаксический разбор, построение абстрактного
синтаксического дерева, разрешение символов, проверку типов и элаборацию (раскрытие параметров,
разворачивание иерархии модулей). Важной особенностью slang является способность обрабатывать
реальный производственный код с поддержкой большинства конструкций стандарта SystemVerilog IEEE
1800-2017.

Внутреннее представление slang включает абстрактное синтаксическое дерево с подробной информацией
о каждом узле, позицию каждого символа в исходных файлах (через объект SourceLocation),
информацию о типах выражений и таблицу символов. Для задачи статического анализа конечных
автоматов особенно важна возможность получать для каждого выражения точную ссылку на его
местоположение в исходном коде --- эта информация впоследствии используется при сопоставлении
результатов анализа с исходным кодом.

\subsection{Синтезатор yosys и его подход к извлечению конечных автоматов}

\subsubsection{Общая информация о Yosys}

Yosys [5] --- открытая инфраструктура для синтеза цифровых логических схем, предназначенная для
преобразования RTL-описаний на Verilog или SystemVerilog в различные представления: от
высокоуровневого списка соединений до вентильного уровня и файлов в форматах, поддерживаемых
конкретными технологиями. Yosys реализует процесс синтеза в виде последовательности независимых
проходов (passes), каждый из которых выполняет определённую трансформацию внутреннего
представления. Такой архитектурный подход обеспечивает расширяемость: новые алгоритмы анализа
и оптимизации могут быть добавлены как отдельные проходы, не затрагивая существующий код.

Внутреннее представление Yosys, называемое RTLIL (RTL Intermediate Language), представляет собой
список логических элементов: проводов (wires), регистров (в том числе D-триггеров различных
типов), мультиплексоров, логических вентилей, арифметических блоков и отдельных специальных
элементов. RTLIL достаточно низкоуровневое, чтобы допускать большое количество оптимизаций,
но при этом достаточно абстрактное, чтобы сохранять высокоуровневые конструкции
(например, ячейка \$add или \$mux не обязательно разворачиваются сразу в вентили).

\subsubsection{Представление логики переходов в виде мультиплексорных деревьев}

Одним из ключевых аспектов работы синтезатора Yosys является приведение различных конструкций языка
SystemVerilog к унифицированному низкоуровневому представлению. За это преобразование в предложенной
схеме анализа отвечает библиотека yosys-slang.

Рассмотрим, как условные конструкции (например, \texttt{if/else} и \texttt{case}) преобразуются в сеть
мультиплексоров. В общем случае любая конструкция ветвления может быть сведена к мультиплексору.
Например, оператор:

\begin{lstlisting}
if (cond)
    a = b;
else
    a = c;
\end{lstlisting}

эквивалентен мультиплексору \texttt{a = cond ? b : c}. Аналогично, оператор case преобразуется в каскад
мультиплексоров с несколькими входами:

\begin{lstlisting}
always_ff @(posedge clk) begin
    case (state)
        S0: state <= S1;
        S1: state <= S2;
        default: state <= S0;
    endcase
end
\end{lstlisting}

На уровне RTLIL данная логика преобразуется в комбинационную схему, формирующую значение следующего
состояния на входе D-триггера. Эта схема представляется в виде дерева мультиплексоров, где:

\begin{itemize}
\item Входы мультиплексоров соответствуют возможным значениям следующего состояния;
\item Управляющие сигналы кодируют текущие условия (например, равенство \texttt{state == S0});
\item Структура дерева отражает вложенность условных операторов.
\end{itemize}

Таким образом, функция переходов конечного автомата реализуется в виде комбинационной схемы, представленной
как мультиплексорное дерево, подающееся на вход D-триггера, хранящего состояние.

\subsubsection{Извлечение конечных автоматов в Yosys}

Для задачи обнаружения и извлечения конечных автоматов в Yosys реализованы два прохода [17]:

\textbf{fsm\_detect} — проход поиска кандидатов в конечные автоматы. Проход перебирает все
D-триггеры в схеме и для каждого проверяет пару (триггер, выходной провод Q). Начиная с D-входа
триггера, выполняется обход мультиплексорного дерева: если на всех его листьях находятся либо
константы, либо тот же самый провод Q, то триггер считается кандидатом в регистр состояния.
Выполняется ряд дополнительных проверок; например, в исходной реализации Yosys требуется, чтобы
выходы автомата зависели только от текущего состояния, что ограничивает обнаружение автоматами Мура.

\textbf{fsm\_extract} — проход собственно извлечения автомата. Он преобразует пару «триггер +
мультиплексорное дерево» в специализированную ячейку \$fsm, инкапсулирующую всю логику автомата:
состояния, таблицу переходов, информацию о сигналах тактирования и сброса [18]. Процедура включает
два этапа: (1) определение входов автомата и возможных состояний обходом мультиплексорного дерева;
(2) определение функций переходов и выходов путём «вычисления» схемы для каждого возможного
значения состояния. Если вычисления недостаточно для однозначного определения следующего состояния,
индивидуальные входы фиксируются в логический ноль или единицу, и процесс повторяется.

Хотя базовые проходы Yosys работоспособны и неплохо покрывают канонические стили конечных
автоматов, они имеют ряд ограничений, существенно сказывающихся на реальных проектах.
Во-первых, не поддерживаются автоматы Мили. Во-вторых, неточно обрабатывается синхронный сброс:
для синхронно сбрасываемых регистров сигнал сброса попадает в число входов автомата, что приводит
к появлению фиктивных переходов из каждого состояния в начальное и искажает последующий анализ
дедлоков и лайвлоков. В-третьих, отсутствует поддержка случаев, когда регистр состояния встроен
в более крупную структуру и на провод отображается не целиком.

В настоящей работе базовая схема (fsm\_detect + fsm\_extract) была сохранена, а каждый проход
индивидуально расширен для устранения описанных ограничений. Детали расширений рассматриваются
в разделе 5.2.

\subsubsection{Ограничения представления}

Несмотря на универсальность, мультиплексорное представление имеет ряд особенностей:

\begin{itemize}
\item Оптимизации синтеза могут изменять структуру дерева (например, объединять ветви или удалять
недостижимые пути);
\item Условия переходов теряют явную семантику и должны восстанавливаться из логической структуры;
\item При наличии сложной комбинационной логики глубина дерева может значительно увеличиваться.
\end{itemize}

Тем не менее, именно такое представление делает возможным унифицированный алгоритм извлечения конечных
автоматов, реализованный в данной работе.

\subsection{Библиотека yosys-slang}

Связующим звеном между библиотекой slang и синтезатором Yosys является библиотека yosys-slang
[16] --- фронтенд, преобразующий АСД, построенное slang, во внутреннее представление RTLIL Yosys.
Такая архитектура имеет несколько преимуществ по сравнению со встроенным в Yosys стандартным
фронтендом Verilog.

Во-первых, slang обеспечивает существенно более точный разбор \\ SystemVerilog, включая расширенные
конструкции стандарта IEEE 1800-2017, которые не поддерживаются штатным фронтендом Yosys. Это
расширяет спектр RTL-описаний, которые может обрабатывать конвейер обработки.

Во-вторых, yosys-slang обеспечивает непрерывную связь между АСД-уровнем (slang) и netlist-уровнем
(RTLIL). Каждый элемент netlist может нести информацию о том, из какой конструкции исходного кода
он был получен. Эта информация хранится в атрибутах RTLIL-ячеек и проводов. В дальнейшем
аттрибуты будут использоваться в работе для восстановления контекста (имён состояний, выражений
условий) при выводе результатов пользователю.

В рамках настоящей работы yosys-slang потребовал расширения: был изменён формат атрибута src,
хранящего информацию о местоположении элемента в исходном коде. Подробности этих изменений
изложены в разделе 5.4.
